\manualchapter{Panel and placement}{sec:panel}
\begin{minipage}[t]{.37\textwidth}
\vspace{0pt}\resizebox{\textwidth}{!}{% Spatial guide to InputGenieWidget in src/CVGenie.cpp.
\begin{tikzpicture}[x=1pt,y=-1pt]
  \draw[panel,rounded corners=3pt] (0,0) rectangle (180,380);
  \draw[panel] (9.5,13) rectangle (170,36);
  \node[label] at (90,25) {GAME SELECTION};
  \PanelBadge{-13}{25}{1}
  \foreach \i in {0,...,7} {
    \pgfmathsetmacro{\y}{69+38.5*\i}
    \PanelJack{26}{\y}
    \draw[detail] (66,\y-12) rectangle (168,\y+12);
    \node[label] at (117,\y) {ROW \the\numexpr\i+1\relax};
  }
  \draw[draw=manualAccent] (-4,55)--(-10,55)--(-10,351)--(-4,351);
  \PanelBadge{-13}{201}{2}
  \node[label] at (90,368) {CV GENIE};
\end{tikzpicture}
}
\end{minipage}\hfill
\begin{minipage}[t]{.58\textwidth}
\setlength{\parskip}{6pt}
\vspace{0pt}
\subsection{1. Game selection}
Click the top display to open \textbf{Games} and choose the running game.
The display initially reads \textbf{No Game Selected}.
Choose a game here to populate the row menus.

Selecting a different game clears the row assignments. It does not change
the loaded ROM; select the matching game in RackNES too.

\subsection{2. Eight input rows}
Each jack controls the selection shown directly to its right. Click the
row display to open \textbf{Game Element}; the menu lists the current map's
entries plus \textbf{Unassigned}. Each entry shows its byte endpoints and
continuous or trigger-toggle mode. Hover over a row display or jack for the
full assignment name and voltage behavior. Rows have no built-in attenuator,
offset, voltage meter, or readback of the game's current value.

Use external CV utilities to scale and offset a source. Leave unused rows
Unassigned. The eight rows are equivalent; the top is row 1, the bottom row 8.
\end{minipage}

\subsection{One adjacent expander}
Only the Genie touching RackNES's right edge communicates with it. A Genie
to the left, on another rack row, or behind a second Genie has no direct
connection. There is no cable for the expander link and no link-status light.
The module has eight monophonic inputs; each reads only the first channel
of a polyphonic cable.

\subsection{Panel themes}
CV Genie and RackNES follow Rack's \textbf{View > Use dark panels if available}
setting (Rack 2.4 or newer). Existing panels and module-browser previews
update together; turn the setting off for light panels. This replaces the
\textbf{Plugin Theme} menu and the legacy \texttt{RackNES.json} preference.
Rack stores the global preference separately from patches.
